% Source: https://artofproblemsolving.com/community/c3h2360922_parabola_question
% Stable identity: aops-topic-2360922-q01
A point $(x,y)$ is on the parabola when its distance to the focus equals its distance to the directrix:
\[\sqrt{(x-1)^2+(y-1)^2}=\dfrac{|x+y|}{\sqrt2}.\]
Squaring and multiplying by $2$:
\[2(x-1)^2+2(y-1)^2=(x+y)^2\iff x^2-2xy+y^2-4x-4y+4=0\iff\boxed{(x-y)^2=4(x+y-1)}.\]
The axis is the line through the focus perpendicular to the directrix, which is the line $\boxed{y=x}$, and the vertex is midway between $F$ and the foot $(0,0)$ of the perpendicular from $F$ to the directrix: $V=\left(\dfrac12,\dfrac12\right)$ (indeed $(0)^2=4(1-1)$).

\begin{asy}
usepackage("amsmath");
size(6cm);
defaultpen(fontsize(10pt));
import graph;
pair Fp=(1,1);
pair f(real t){ real u=(t^2+2)/(2*sqrt(2)); return ((u+t)/sqrt(2),(u-t)/sqrt(2)); }
draw((-1.6,0)--(4.2,0),Arrow(5)); draw((0,-1.6)--(0,4.2),Arrow(5));
draw(graph(f,-2.5,2.5),blue+linewidth(0.9));
draw((-1.5,1.5)--(1.5,-1.5),red+dashed); label("$x+y=0$",(-1.5,1.5),W,red+fontsize(8pt));
draw((-0.8,-0.8)--(3.6,3.6),dotted); label("axis $y=x$",(3.6,3.6),NE,fontsize(8pt));
dot(Fp); label("$F(1,1)$",Fp,NE,fontsize(8pt)); dot((0.5,0.5)); label("$V$",(0.5,0.5),1.5*dir(-80),fontsize(8pt));
\end{asy}

So \textbf{yes}: this parabola's axis, $y=x$, is parallel to neither coordinate axis. (In general, a parabola can have any line as its axis; its equation then contains an $xy$ term.)
